\section{Conclusiones}

La Etapa 1 terminó con tres buenos baselines y una pregunta incómoda: ¿de qué
sirve superar 0.90 de macro-F1 si potasio sigue fallando y el test ya fue visto
varias veces? La segunda etapa no resolvió ese problema con una sola cifra.
Cambió el corpus, congeló un holdout nuevo, comparó representaciones bajo el
mismo protocolo y mantuvo las métricas minoritarias junto al resultado global.

El tuning produjo la mejora controlada más clara del trabajo. Sobre
EfficientNet-B0 y la misma validación, la cabeza pasó de 0.7799 a 0.8887 de
macro-F1. El cambio que más importó no fue añadir parámetros: fue reducir la
tasa, usar penalización L2 y ponderar las clases con la raíz de su frecuencia
inversa. Recall macro aumentó 0.1399. Esto confirma que el baseline dejaba
información útil en los embeddings, pero la frontera aprendida favorecía a las
clases grandes.

EfficientNet-B0 volvió a ser el mejor modelo individual, aunque no dominó cada
imagen. Su conjunto de errores tuvo Jaccard 0.277 con MobileNetV3; esa
complementariedad permitió que el voto ponderado alcanzara 0.9202 en
validación, 0.0315 más que B0. El costo es concreto: cuatro backbones, cerca de
38.3 MB FP32 sumados y más latencia. Para investigación se selecciona el
ensemble; para teléfono, ShuffleNet o MobileNet siguen siendo alternativas que
merecen una prueba de dispositivo.

Los cinco folds situaron el macro-F1 en $0.9097\pm0.0090$. El holdout único
quedó en 0.9035, acompañado por 0.9537 de accuracy. La cercanía con validación
cruzada indica que la selección no dependió completamente de un corte, pero el
quinto fold reveló su límite: potasio cayó a F1 0.6304. En el holdout, potasio
obtuvo 0.6706 y solo 57 de 93 imágenes correctas. Veintiocho de sus 36 errores
fueron nitrógeno. La clase mejoró frente a la referencia histórica de 0.62,
pero sigue siendo el problema central y no debe ocultarse detrás de 95\% de
accuracy.

El sesgo laboratorio/campo resultó más específico de lo que sugería un promedio
global. Mancha gris y tizón tuvieron recalls ligeramente mayores en campo;
roya pasó de 1.0 en 322 imágenes de laboratorio a 0.3125 en 16 de campo. El
soporte pequeño impide fijar la magnitud real, pero la dirección confirma que
el fondo controlado puede sostener un atajo. Las brechas por fuente y clase
fueron mayores que las de resolución. La acción correcta no es declarar que el
modelo es justo, sino recolectar roya de campo, revisar N/P/K y monitorear cada
segmento con su soporte.

La aplicación móvil completa la cadena técnica: captura, preprocesamiento,
TFLite local, confianza, OOD, resultado e historial offline. Sus 276 pruebas y
typecheck aprobaron; la demo experimental también produjo una salida real. Sin
embargo, el modelo final de esta etapa no reemplazó el EfficientNet-Lite0
embarcado y el despliegue público no estuvo accesible al cierre. El informe
descuenta esa evidencia faltante. Un prototipo funcional no se convierte en
servicio disponible solamente porque exista una URL en la configuración.

\subsection{Cumplimiento de la rúbrica}

La calificación documentada es 96/100. Los cuatro puntos no otorgados
corresponden al criterio de prototipo: faltó comprobar disponibilidad pública
continua y ejecutar el modelo experimental final dentro de Android. Los demás
criterios poseen CSV, JSON, código, figuras y secciones enlazadas. No se asignó
puntaje por una capacidad registrada sin corrida.

Doctor Maíz queda en un punto útil pero delimitado. Ya existe evidencia
reproducible de selección y una aplicación donde integrar modelos; todavía
falta la evidencia que solo puede dar el campo. La siguiente decisión no
debería buscar otra arquitectura por novedad, sino comprobar con productores y
extensionistas si las hojas locales, los teléfonos reales y la forma de mostrar
incertidumbre sostienen lo observado aquí.

\begin{table}[H]
  \centering
  \caption{Puntuación final respaldada por artefactos reproducibles.}
  \label{tab:rubric}
  \small
  \begin{tabularx}{\textwidth}{>{\raggedright\arraybackslash}X r l >{\raggedright\arraybackslash}X}
\toprule
\textbf{Criterio} & \textbf{Puntos} & \textbf{Estado} & \textbf{Evidencia} \\
\midrule
Optimización e hiperparámetros & 20/20 & Completo & 30 trials, 13 podados; macro-F1 0.7799 a 0.8887. \\
Modelos avanzados y ensemble & 20/20 & Completo & Cuatro backbones; weighted soft voting mejora 0.0315. \\
Evaluación rigurosa & 15/15 & Completo & 5-fold y holdout congelado evaluado una sola vez. \\
Sesgos y ética & 15/15 & Completo & Brechas por clase, entorno, fuente, blur y resolución. \\
Prototipo & 11/15 & Parcial & Demo y 276 pruebas; API/APK no disponibles al cierre. \\
Impacto y recomendaciones & 15/15 & Completo & Condiciones de uso y nueve propuestas concretas. \\
\midrule
\textbf{TOTAL} & \textbf{96/100} & Cumplimiento con una salvedad & La disponibilidad pública del despliegue no pudo confirmarse al cierre. \\
\bottomrule
\end{tabularx}

\end{table}
